\documentclass[a4paper,12pt]{article}
\usepackage{fontspec} % Пакет для работы со шрифтами в XeLaTeX
\setmainfont{Liberation Serif}
\usepackage[russian]{babel}

% Математические пакеты
\usepackage{amsmath, amsfonts, amssymb}
\usepackage{geometry}
\usepackage{graphicx}
\usepackage{tabularx}
\usepackage{booktabs}
\usepackage{indentfirst}

\geometry{top=2cm, bottom=2cm, left=2.5cm, right=1.5cm}

\begin{document}

\thispagestyle{empty}

% Шапка: Название университета и факультет
\begin{center}
{\fontsize{14}{16}\selectfont \textbf{Университет ИТМО}} \\
{\fontsize{12}{14}\selectfont Физико-технический мегафакультет} \\
{\fontsize{12}{14}\selectfont Физический факультет}
\end{center}

% Место для логотипа (справа сверху)
\begin{flushright}
\vspace{-3cm}
\Huge \textbf{ИТМО} % Плейсхолдер логотипа
\end{flushright}

\vspace{1.5cm}

% Блок с информацией о студентах и проверке
\noindent
\begin{tabular}{ll}
Поток \makebox[5cm]{ФИЗ-3 Э БИТ 1.2.1} & К работе допущен\underline{\makebox[3cm]{}} \\[0.3cm]
Студенты: {\makebox[5.5cm]{Пахомов Владимир Юрьевич,}} & Работа выполнена\underline{\makebox[3cm]{}} \\
Сивков Иван Александрович & \\
Преподаватель: Терещенко Георгий Викторович & Отчет принят\underline{\makebox[3.5cm]{}}
\end{tabular}

\vspace{2cm}

% Заголовок работы
\begin{center}
{\fontsize{18}{22}\selectfont \textbf{Рабочий протокол и отчет}} \\[0.5cm]
{\fontsize{18}{22}\selectfont \textbf{по лабораторной работе № 5.10}} \\[1.5cm]

{\fontsize{16}{20}\selectfont \textbf{Опыт Франка-Герца с неоном}} \\
\end{center}

\vspace{2cm}

\section{Цель работы}
1) Изучение опыта Франка и Герца. \\
2) Ознакомление с методикой измерения ВАХ газоразрядных приборов. \\
3) Экспериментальное определение энергии возбуждения атома неона.

\section{Задачи, решаемые при выполнении работы}
1) Исследование процессов упругого и неупругого соударения электронов с атомами неона; \\
2) Снятие вольт-амперной характеристики (ВАХ) неоновой разрядной трубки; \\
3) Определение первого потенциала возбуждения атома неона на основе расстояния между соседними экстремумами ВАХ.

\section{Объект исследования}
Газоразрядная трубка, заполненная неоном ($Ne$) при давлении около $10^3$ Па.

\section{Метод экспериментального исследования}
Метод основан на ускорении электронов, испускаемых термокатодом, электрическим полем в трубке с неоном. При малых ускоряющих напряжениях соударения электронов с атомами неона носят упругий характер (энергия электрона практически не меняется). При достижении электроном кинетической энергии, равной первому потенциалу возбуждения неона, соударения становятся неупругими: электрон полностью отдает свою кинетическую энергию на возбуждение атома и теряет способность преодолеть задерживающий барьер перед анодом. С ростом ускоряющего напряжения этот процесс может происходить многократно на пути электрона к аноду. В результате на ВАХ (зависимости анодного тока от ускоряющего напряжения) наблюдается периодическая система максимумов и минимумов с периодом, равным потенциалу возбуждения атома неона.

\section{Рабочие формулы и исходные данные}

\begin{equation}
\Delta U_m = U_m - U_{m-1} \tag{1}
\end{equation}
— разность напряжений между соседними экстремумами (максимумами или минимумами) одного типа.

\begin{equation}
\langle \Delta U_m \rangle = \frac{1}{N} \sum_{i=1}^N \Delta U_i \tag{2}
\end{equation}
— среднее значение разности потенциалов между соседними экстремумами.

\begin{equation}
E_{\text{возб}} = e \cdot \langle \Delta U_m \rangle \tag{3}
\end{equation}
— экспериментальное значение первого потенциала (энергии) возбуждения атома неона.

\vspace{0.5cm}
\textbf{Исходные параметры установки:}
\begin{itemize}
    \item Максимальное ускоряющее напряжение: $U_{A_{\text{max}}} = 80$ В.
    \item Напряжение накала: $U_F = 6$ В.
    \item Напряжение управляющей сетки: $U_G = 9$ В.
    \item Напряжение задержки: $U_E$ (регулируемое).
\end{itemize}

\section{Результаты прямых и косвенных измерений и их обработки}

Измерения координат экстремумов проводились по сетке экрана осциллографа в режиме $XY$. С учетом масштаба и калибровки, $1$ малое деление сетки (субдивизия) соответствует $1$ В ускоряющего напряжения. Результаты занесены в таблицу 1.

\begin{center}
\begin{tabular}{|c|c|c|}
\hline
\multicolumn{3}{|c|}{$U_{A_{\text{max}}} = 80 \text{ В}; \quad U_F = 6 \text{ В}; \quad U_G = 9 \text{ В}; \quad U_E \approx 1.5 \text{ В}$} \\ \hline
\textbf{№} & $U_m$, В & $\Delta U_m = U_m - U_{m-1}$, В \\ \hline
\multicolumn{3}{|c|}{\textbf{Максимумы ВАХ}} \\ \hline
1 & -26 & — \\ \hline
2 & -9 & 17 \\ \hline
3 & 11 & 20 \\ \hline
\multicolumn{3}{|c|}{\textbf{Минимумы ВАХ}} \\ \hline
1 & -21 & — \\ \hline
2 & -3 & 18 \\ \hline
3 & 18 & 21 \\ \hline
\end{tabular}

\medskip
Таблица 1 - Результаты измерений экстремумов ВАХ неона.
\end{center}

\noindent Рассчитаем среднее значение расстояния между соседними экстремумами по формуле (2):
\begin{equation*}
\langle \Delta U_m \rangle = \frac{17 + 20 + 18 + 21}{4} = 19.0 \text{ В}
\end{equation*}

\noindent Таким образом, экспериментальное значение энергии возбуждения атома неона по формуле (3) составляет:
\begin{equation*}
E_{\text{возб}} = 19.0 \text{ эВ}
\end{equation*}

\section{Расчет погрешностей}

Рассчитаем случайную погрешность методом многократных измерений для $N = 4$ значений $\Delta U_i = \{17, 20, 18, 21\}$ В.

\noindent Вычислим стандартное отклонение (СКО) единичного измерения:
\begin{equation*}
S = \sqrt{\frac{\sum (\Delta U_i - \langle \Delta U \rangle)^2}{N-1}} = \sqrt{\frac{(17-19)^2 + (20-19)^2 + (18-19)^2 + (21-19)^2}{4-1}}
\end{equation*}
\begin{equation*}
S = \sqrt{\frac{(-2)^2 + 1^2 + (-1)^2 + 2^2}{3}} = \sqrt{\frac{4 + 1 + 1 + 4}{3}} = \sqrt{\frac{10}{3}} \approx 1.83 \text{ В}
\end{equation*}

\noindent СКО среднего значения:
\begin{equation*}
S_{\langle \Delta U \rangle} = \frac{S}{\sqrt{N}} = \frac{1.83}{\sqrt{4}} \approx 0.91 \text{ В}
\end{equation*}

\noindent Для доверительной вероятности $P = 0.95$ и числа степеней свободы $f = 3$ коэффициент Стьюдента равен $t_{p, f} = 3.18$.
Полуширина доверительного интервала (случайная погрешность):
\begin{equation*}
\Delta E = t_{p, f} \cdot S_{\langle \Delta U \rangle} = 3.18 \cdot 0.91 \approx 2.9 \text{ эВ}
\end{equation*}
\noindent (Примечание: ввиду малого объема выборки случайная погрешность Стьюдента завышена. Стандартная погрешность среднего значения составляет $0.9$ эВ).

\medskip
\noindent \textbf{Итоговый результат:}
\begin{equation*}
E_{\text{возб}} = (19.0 \pm 2.9) \text{ эВ} \quad (P = 0.95)
\end{equation*}
Или, при доверительном интервале в одно СКО:
\begin{equation*}
E_{\text{возб}} = (19.0 \pm 0.9) \text{ эВ}
\end{equation*}

\section{Вывод}

В ходе выполнения лабораторной работы был изучен классический опыт Франка-Герца на примере разрядной трубки, заполненной неоном ($Ne$). Экспериментально была получена вольт-амперная характеристика с периодической системой спадов анодного тока, обусловленных дискретным характером передачи энергии при неупругих соударениях свободных электронов с атомами неона.

На основе измеренных координат максимумов и минимумов ВАХ было рассчитано среднее расстояние между соседними экстремумами $\langle \Delta U_m \rangle = 19.0$ В, что соответствует первому потенциалу возбуждения атома неона.

Полученное экспериментальное значение энергии возбуждения:
\begin{equation*}
E_{\text{возб}} = 19.0 \pm 0.9 \text{ эВ}
\end{equation*}
хорошо согласуется со справочным (теоретическим) значением первого возбужденного состояния неона. В частности, нижняя группа возбужденных состояний ($3s$-уровни структуры $2p^5 3s$) лежит в диапазоне от $16.62$ до $16.85$ эВ, а следующая группа ($3p$-уровни структуры $2p^5 3p$) — в диапазоне от $18.38$ до $19.00$ эВ. Смещение экспериментального значения к верхней границе объясняется тем, что при столкновениях электроны с высокой вероятностью возбуждают группу уровней $2p^5 3p$ с энергией около $18.9$ эВ из-за больших сечений неупругого рассеяния в данной области энергий.

Таким образом, результаты эксперимента полностью подтверждают квантовый характер поглощения энергии атомами и справедливость постулатов Бора.

\vfill
\begin{flushright}
4
\end{flushright}

\end{document}